% ─── Diagramme de classes — Entités principales ──────────────────────────────
% La macro \umlclass est définie dans le préambule de main.tex
% Layout : 3 colonnes × 2 rangées, espacées de 6cm horizontalement
%          et 7.5cm verticalement pour éviter tout chevauchement

\begin{tikzpicture}[
  inh/.style={-{Triangle[open, length=7pt]}, thick, draw=primary},
  assoc/.style={-{Stealth[length=5pt]}, draw=gray!70, thin},
  comp/.style={-{Diamond[open, length=7pt]}, draw=gray!70, thin},
  lbl/.style={font=\scriptsize, fill=white, inner sep=1pt, midway},
]

% ── Rangée haute : User  |  StudentProfile  |  RecruiterProfile ──────────────
\umlclass{user}{0}{0}{4.2cm}{User}{%
\textbf{PK} id : UUID\\
email : String\\
username : String\\
password : String\\
role : Role}{%
+register()\\
+login()}

\umlclass{student}{6}{0}{4.2cm}{StudentProfile}{%
\textbf{FK} userId : UUID\\
bio : String?\\
currentYear : Int?\\
preferredDomains : String[\,]\\
preferredWorkMode : String[\,]\\
preferredCities : String[\,]}{%
+update()\\
+uploadAvatar()}

\umlclass{recruiter}{12}{0}{4.2cm}{RecruiterProfile}{%
\textbf{FK} userId : UUID\\
company : String\\
website : String?\\
logoUrl : String?\\
description : String?}{%
+update()}

% ── Rangée basse : Offer  |  Application  |  Notification ────────────────────
\umlclass{offer}{0}{-8}{4.2cm}{Offer}{%
\textbf{PK} id : UUID\\
\textbf{FK} recruiterId : UUID\\
title : String\\
description : String\\
type : OfferType\\
workMode : WorkMode\\
isPaid : Boolean\\
stipendMin : Float?\\
stipendMax : Float?}{%
+create()\\
+update()\\
+delete()}

\umlclass{application}{6}{-8}{4.2cm}{Application}{%
\textbf{PK} id : UUID\\
\textbf{FK} studentId : UUID\\
\textbf{FK} offerId : UUID\\
status : AppStatus\\
cvId : UUID?\\
coverLetterId : UUID?\\
createdAt : DateTime}{%
+submit()\\
+updateStatus()\\
+withdraw()}

\umlclass{notif}{12}{-8}{4.2cm}{Notification}{%
\textbf{PK} id : UUID\\
\textbf{FK} userId : UUID\\
type : String\\
title : String\\
message : String\\
read : Boolean\\
createdAt : DateTime}{%
+markRead()\\
+markAllRead()}

% ── Relations ─────────────────────────────────────────────────────────────────
% User ←extends— StudentProfile (héritage/composition 1-1)
\draw[inh] (student_title.north) -- ++(0,0.5)
  -| node[lbl, near start]{1} (user_title.north)
  node[lbl, pos=0.85, right]{\scriptsize «FK»};

% User ←extends— RecruiterProfile
\draw[inh] (recruiter_title.north) -- ++(0,0.5)
  -| node[lbl, near start]{1} (user_title.north east);

% RecruiterProfile 1 ——< N Offer
\draw[comp] (offer.north) -- ++(0,0.7) -|
  node[lbl, near start]{N}
  node[lbl, near end, right]{1}
  (recruiter.south);

% StudentProfile 1 ——< N Application
\draw[comp] (application.north) -- ++(0,0.7) -|
  node[lbl, near start]{N}
  node[lbl, near end, right]{1}
  (student.south);

% Offer 1 ——< N Application
\draw[assoc] (application_title.west) -- node[lbl, above]{N\quad 1} (offer_title.east);

% User 1 ——< N Notification
\draw[assoc] (notif_title.north) -- ++(0,0.8) -|
  node[lbl, near start]{N}
  node[lbl, near end, right]{1}
  (user_title.south);

% ── Énumérations — placées à l'extérieur de la grille de classes ──────────────
% Role : à gauche de User (row 1), hors de la grille
\node[draw=orange!60, fill=orange!8, rounded corners=3pt,
      font=\scriptsize, align=left, inner sep=5pt] (enum_role) at (-4.8, 0) {
  \textbf{\guillemotleft enum\guillemotright}\\
  \textbf{Role}\\
  STUDENT\\RECRUITER\\ADMIN
};

% OfferType : à gauche de Offer (row 2), hors de la grille
\node[draw=orange!60, fill=orange!8, rounded corners=3pt,
      font=\scriptsize, align=left, inner sep=5pt] (enum_type) at (-4.8, -8) {
  \textbf{\guillemotleft enum\guillemotright}\\
  \textbf{OfferType}\\
  INTERNSHIP\\PFE\\RESEARCH\\PHD\\ALTERNANCE
};

% ApplicationStatus : à droite de Application (row 2), hors de la grille
\node[draw=orange!60, fill=orange!8, rounded corners=3pt,
      font=\scriptsize, align=left, inner sep=5pt] (enum_status) at (16.0, -8) {
  \textbf{\guillemotleft enum\guillemotright}\\
  \textbf{ApplicationStatus}\\
  SUBMITTED\\IN\_REVIEW\\ACCEPTED\\REJECTED\\WITHDRAWN
};

\end{tikzpicture}
